\centerline{\bf \Huge Разминочная перед графами}
\title{Разминочная перед графами}

\begin{flushright}
    \scriptsize Реро реро реро реро реро реро реро реро...\\Нориаки Какёин
\end{flushright}

\textit{Формула суммы арифметической прогрессии:} $S_n=\dfrac{(a_1 + a_n)d}{2} = \dfrac{2a_1 + d(n-1)}{2}\cdot n$

\textit{Формула суммы геометрической прогрессии:} $\dfrac{b_1 (q^n-1)}{q-1}$
    
\problem{1} а)Алине за летние сборы надо решить  742 задачи. Ежедневно она решает на одно и то же количество задач больше по сравнению с предыдущим днем. Известно, что за первый день Алина решила 5 задач. Определите, сколько задач решила Алина в последний день, если со всеми задачами она справилась за 14 дней.\\
б)Бизнесмен Доржеев получил в 2023 году прибыль в размере 50000 рублей. Каждый следующий год его прибыль увеличивалась на 300\% по сравнению с предыдущим годом. Сколько рублей заработал Доржеев за 2027 год?

\problem{2}Найдите сумму:\\
а)$1+2+3\dots+(2n-1)=?$\\
б)$1+2+4+8+\dots+2^n=?$

\problem{3}Найти сумму:\\
а) 1+11+111+...+111...1, где последнее число содержит n единиц; \\
б)аналогичная задача, когда вместо единиц стоят пятерки.

\problem{4}Герман своровал в полдник 777 чокопаев. Герман хочет съесть все чокопаи за n дней, причем так, чтобы каждый из этих дней (кроме первого, но включая последний) съедать на один чокопай больше, чем в предыдущий. Для какого наибольшего числа n это возможно?

\problem{5}
Баир изобрёл машину времени и решил навестить маленького Гаусса, далеко в прошлом. Гаусс решил показать только что изобретённый способ складывать числа от 1 до какого-то n. И для примера Гаусс взял какое-то интересное число. Так как Баир очень долго создавал машину времени, в том числе очень долго её программировал, посмотрев на полученную сумму он сказал:" А в двоичном коде это число(сумма чисел) выглядело как 10 подряд единичек, а все остальные нули. Действительно ли есть такое интересное число, или Баир из-за усталости немного просчитался?

\newpage

\problem{6}Можно ли взять несколько арифметических прогрессий так, чтобы совокупности в них находились все натуральные числа?

\problem{7}Компания ЛмШ решила телепортировать НеНачальника прямиком в ловушку. И что бы он не сумел бежать по краям телепортационного круга решили устроить формирование из Лососей.

1) По одному Лососю поставили на концы диаметра круга.

2) По среди всех новых дуг образованными лососями ставили еще Лососей. Причем в каждую дугу ставили столько Лососей же сколько суммарно Лососей на конце новой дуги.

Пункт 2) повторили n раз. Сколько Лососей участвовали в засаде на НеНачальника?

\problem{8}Дана арифметическая прогрессия, которая начинается с 1. За каждую найденную геометрическую прогрессию длины n НеНачальник обещает дать 1000001 биткоинов. Но если НеНачальник сможешь найти прогрессию большей длины, то ничего не даст. НеНачальник опять всех обманывает?

\vspace{\baselineskip}

\problem{9}.\textbf{\Huge УЛЬТРА МЕГА ГРОБ ОТ РАСУЛА.}Назовём \textit{паспортом} натурального числа набор простых чисел, входящих в его разложение. Докажите, что в любой натуральной арифметической прогресии есть сколько угодно много чисел с одинаковым паспортом.